\documentclass[12pt,a4paper]{article}
\usepackage[utf8]{inputenc}
\usepackage[T1]{fontenc}
\usepackage{amsmath}
\usepackage{hyperref}
\usepackage{geometry}
\usepackage{enumitem}
\usepackage{tabularx}
\usepackage{booktabs}
\usepackage{parskip}
\geometry{margin=2.5cm}
\setlength{\parindent}{0pt}

\title{Cancer Screening Microsimulation Model\\--- Plain-Language Description ---}
\author{}
\date{\today}

\begin{document}
\maketitle

\section{What This Model Is}

This is an agent-based microsimulation: a computer program that creates a large virtual population and tracks each person individually from birth to death. Every person can develop cancer, be diagnosed, receive treatment, participate in screening, and die---either from cancer or from other causes. By comparing outcomes with and without screening, the model estimates how many lives a screening programme can save, and how many years of life it adds.

The model uses real demographic and epidemiological data to calibrate its parameters, so that the simulated incidence and mortality curves match those observed in the population.

\section{What Each Virtual Person Looks Like}

Each person is represented by a set of numeric attributes stored as columns in a table (one row per person):

\begin{itemize}[nosep]
    \item \textbf{Year of birth} --- derived from the population age distribution.
    \item \textbf{Natural death age} --- sampled from the all-cause mortality distribution.
    \item \textbf{Cancer diagnosis age} --- the age at which cancer would be diagnosed clinically, without screening. Generated from the diagnosis hazard.
    \item \textbf{Has cancer} --- true if the diagnosis age is earlier than the natural death age.
    \item \textbf{Stage at diagnosis} --- I, II, III, or IV. The number of stage-size thresholds the tumour has crossed by the diagnosis time (Gompertz).
    \item \textbf{Aggressiveness} --- whether the tumour is aggressive; a cure-only label drawn per stage.
    \item \textbf{Tumour growth rate} \; $B$ --- the individual Gompertz rate (how fast the tumour progresses through stages).
    \item \textbf{Incidence age} --- the age when the first malignant cell appears (years before diagnosis).
    \item \textbf{Stage ages} --- the ages at which the tumour reaches each stage (0 through terminal), from the Gompertz crossing times.
    \item \textbf{Cured} --- whether clinical treatment was successful.
    \item \textbf{Screening-detected} --- whether the cancer was found by screening.
    \item \textbf{Screening-cured} --- whether treatment after screening detection was successful.
    \item \textbf{Cancer death ages} --- three variants: natural history (no treatment), with clinical cure, and with screening cure.
\end{itemize}

\section{The Life Story of One Person}

\subsection{Step 1: Initial Age and Natural Death}

Each person starts at a random age drawn from the population age distribution. Their natural death age is sampled from the all-cause mortality distribution (the probability of dying at each age, derived from life tables).

\subsection{Step 2: Will They Get Cancer?}

The probability of being clinically diagnosed at age~$a$ is given by a piecewise-constant function fitted to real incidence data. The diagnosis age is generated by checking at each age whether diagnosis occurs (a Bernoulli trial with age-dependent probability). If the sampled diagnosis age exceeds the natural death age, the person never develops cancer.

\subsection{Step 3: Cancer Natural History}

If diagnosed before natural death, the model reconstructs the unobserved preclinical trajectory:

\begin{itemize}[nosep]
    \item \textbf{Growth rate} is drawn per person: $B = \exp(\log B_{\text{pop}} + \varepsilon)$ with $\varepsilon \sim N(0, \sigma_B^2)$, so people differ in how fast their tumour grows.
    \item \textbf{Stage crossing times} follow from the reduced Gompertz curve $V(t) = \exp(K - \exp(C - B t))$: the tumour reaches size $s$ at $t_s = (C - \log(K - \log s))/B$.
    \item \textbf{Lead time} (time from first malignant cell to clinical diagnosis) is drawn from an exponential distribution with mean \texttt{LeadTimeMean}.
    \item \textbf{Stage at diagnosis} is the number of size thresholds ($2,3,4$) crossed by the lead time.
    \item \textbf{Aggressiveness} is drawn from the observed proportion of aggressive tumours for that stage; it affects \emph{only} the cure odds, not the growth.
\end{itemize}

The per-stage crossing times are \emph{derived from the staging data's stage mix}, so the model reproduces the observed distribution of stages at diagnosis. The only hand-set quantity is the overall time scale (\texttt{LeadTimeMean}); the cancer death hazard is a constant exponential rate $\lambda$.

\subsection{Step 4: Treatment Outcome}

At clinical diagnosis or screening detection, treatment is applied. The cure probability is:

\begin{equation}
    P_{\text{cure}} = \min\left( \; \text{AgeEffect}(a) \times \eta_s, \;\; 1 \right),
\end{equation}

where $\eta_s$ is the stage-specific treatment efficacy, and $\text{AgeEffect}(a)$ is a piecewise-constant multiplier with five age groups (0--40, 40--50, 50--60, 60--70, 70--80, 80+). If cured, the cancer death age is pushed past the natural death age.

\subsection{Step 5: Year-by-Year Simulation}

Each simulation year, the following happens for every alive person:

\begin{enumerate}[nosep]
    \item \textbf{Natural death}: if the person has reached their natural death age, they die.
    \item \textbf{Screening}: if the screening programme is active this year, eligible people may be screened (see below).
    \item \textbf{Cancer death}: if the person has cancer and their cancer death age has arrived, they die.
    \item \textbf{Statistics}: counts of alive people, new cases, and deaths are recorded by age.
\end{enumerate}

\subsection{Step 6: Screening}

Screening operates in certain years, at a fixed interval (e.g.\ every~2 years), for an eligible age range (e.g.\ 50--60). For each alive, undiagnosed person in that range:

\begin{itemize}[nosep]
    \item With probability equal to the participation rate, they are screened.
    \item If they have cancer that started before their current age, the test detects it with probability TP (sensitivity).
    \item If they do not have cancer, the test returns a false positive with probability FP ($1 - \text{specificity}$).
    \item If cancer is found, the current stage is determined, cure is attempted, and the death age is recalculated.
\end{itemize}

A person is screened at most once.

\section{What the Model Calculates}

After the full simulation, the model aggregates:

\begin{itemize}[nosep]
    \item \textbf{Incidence rates} --- new cancer cases per year, by age.
    \item \textbf{Mortality rates} --- cancer deaths per year, by age, both overall and screening-affected.
    \item \textbf{Stage distribution} --- fraction of cancers found at each stage, clinically vs.\ by screening.
    \item \textbf{Survival} --- cause-specific survival function $S(t)$ from diagnosis, estimated as $\exp(-H(t))$ where $H(t)$ is the cumulative hazard.
    \item \textbf{Years of life saved} --- total person-years gained because of screening.
    \item \textbf{People saved} --- number of people whose cancer death was prevented by screening.
\end{itemize}

\section{Model Parameters}

\begin{center}
\begin{tabularx}{\textwidth}{@{} l X @{}}
\toprule
\textbf{Parameter} & \textbf{What it controls} \\
\midrule
Years to simulate & How many years the simulation runs (e.g.\ 30). \\
Initial population & How many virtual people are created (e.g.\ 100\,000). \\
Diagnosis hazard & Age-specific probability of clinical cancer diagnosis. Fitted to real data. \\
Cancer death hazard & Rate at which diagnosed cancer causes death (constant exponential rate). \\
Lead time mean & Mean preclinical duration (the single non-identifiable time scale). The stage crossing times are derived from the staging data. \\
Gompertz parameters & Shape of the tumour growth curve, fitted once on all records (one model for all). \\
Treatment efficiency & Probability of cure for stages\;I--IV. \\
Age cure modifiers & How cure probability changes with age (5 age groups). \\
Recurrence probability & Chance that cancer returns after treatment (flag only; does not affect mortality). \\
Screening start / end age & Who is eligible for screening. \\
Screening interval & How often screening occurs (e.g.\ every 2\;years). \\
Participation rate & Fraction of eligible people who actually get screened. \\
Test TP / FP & Sensitivity and false-positive rate of the screening test. \\
\bottomrule
\end{tabularx}
\end{center}

\section{How the Model is Calibrated}

Calibration adjusts three sets of parameters so that simulated outputs match real population data:

\begin{enumerate}[nosep]
    \item \textbf{Diagnosis hazard.} The age-specific diagnosis probabilities are fitted via Poisson maximum likelihood: \textit{expected cases} = incidence probability $\times$ population at each age. This maximises the match to observed case counts.
    \item \textbf{Tumour growth (Gompertz model).} The stage-crossing times are derived from the observed stage mix, and the growth-curve parameters are fitted to them so the model reproduces the distribution of stages at diagnosis.
    \item \textbf{Cancer death hazard.} The constant death rate~$\lambda$ is fitted analytically via Poisson likelihood, matching observed mortality counts by age.
\end{enumerate}

Calibration uses standard numerical optimisers (L-BFGS-B with bounds for the Gompertz fit) and takes about a second for a typical dataset.

\section{Assumptions (What the Model Simplifies)}

\begin{enumerate}
    \item \textbf{Population homogeneity.} Everyone shares the same parameters except for random variation. No stratification by sex, geography, or socioeconomic status. Risk factors (smoking, obesity) are defined but not yet used.

    \item \textbf{Event independence.} Natural death, cancer diagnosis, treatment outcome, and screening participation are independent processes. Correlations (e.g.\ frailty affecting both mortality and treatment response) are not modelled.

    \item \textbf{No temporal trends.} Parameters are constant over the whole simulation horizon. Improvements in diagnostics, treatment, or screening technology over time are not simulated.

    \item \textbf{Constant cancer death hazard.} The risk of dying from cancer is assumed memoryless (exponential distribution). A time-dependent hazard (Weibull, log-logistic) would be more biologically realistic and is available in the code but not used by default.

    \item \textbf{Deterministic stage transitions (per person).} Given the individual rate $B$, the crossing times are fixed by the Gompertz curve. Randomness in progression comes only from the per-person $B$ (lognormal, spread $\sigma_B$).

    \item \textbf{Complete cure.} A cured person lives to their natural death age. Partial remission, treatment side-effects, and late recurrence are not modelled. The recurrence flag exists but does not affect survival.

    \item \textbf{No pre-cancer stage.} Carcinoma in situ is not included. All cancers arise directly as invasive tumours. This may overstate screening benefit, because some screen-detected pre-cancers would never have progressed.

    \item \textbf{Closed population.} No births, no migration. The initial population only ages and dies.

    \item \textbf{Discrete time.} Time steps are 1\;year. All events occur at year boundaries; within-year dynamics are ignored.

    \item \textbf{One screening event.} After screening detects cancer, that person is excluded from further screening.

    \item \textbf{Lead time is an assumption, not a measurement.} How long a tumour grows before diagnosis cannot be observed directly (diagnosis triggers treatment, interrupting the natural history). The model assumes lead times follow an exponential distribution with user-specified means. Different assumptions produce different results for screening outcomes; the sensitivity analysis tool quantifies this effect.

    \item \textbf{No treatment-related mortality.} Cured patients face no extra risk of death from surgery or chemotherapy complications. The parameter for this exists but is not activated.

    \item \textbf{No behavioural change.} A false-positive screening result does not alter future participation behaviour.

    \item \textbf{Uniform treatment.} All patients at the same stage receive the same chance of cure, regardless of comorbidities or healthcare access.
\end{enumerate}

\section{Web Interface}

The model includes a browser-based interface:

\begin{itemize}[nosep]
    \item \textbf{Fit + Simulate} --- automatically calibrates parameters to real data, then runs the simulation. Use this on first run so results are realistic.
    \item \textbf{Run Only} --- runs the simulation with previously calibrated (or manually edited) parameters.
    \item \textbf{Parameters panel} --- view and edit all model parameters. Changes take effect on the next run.
    \item \textbf{Save / Load} --- save calibrated parameters to a JSON file and reload them later, avoiding repeated calibration.
    \item \textbf{Sensitivity Analysis} --- runs the model 5 times with lead time assumptions multiplied by 0.5, 0.75, 1.0, 1.25, and 1.5, showing overlaid incidence, mortality, and survival curves so you can see how much the conclusions depend on the unverifiable lead-time assumption.
\end{itemize}

\section{Sensitivity Analysis}

The lead time (time from first cancer cell to clinical diagnosis) is the main unobservable quantity in the model. The sensitivity analysis shows what happens if the assumed lead times are half or double the default. The results are shown as multiple curves on the same graphs. If the curves nearly overlap, the model is robust to this uncertainty. If they spread widely apart, the conclusions about screening effectiveness depend critically on the lead-time assumptions and should be treated cautiously.

\section{Running the Model}

\subsection{Command Line}

\begin{verbatim}
cd MedicalModelPython

:: Quick test:
python run_simulation.py --population 50000 --seed 42

:: With calibration:
python run_simulation.py --population 100000 --fit --seed 12345
\end{verbatim}

\subsection{Web Interface}

\begin{verbatim}
python web_app.py
\end{verbatim}

Open \texttt{http://127.0.0.1:5000} in a browser. Click \textbf{Fit + Simulate} first.

\end{document}